\documentclass[UTF8,fontset=fandol,aspectratio=169,11pt]{ctexbeamer}
% Shared Beamer preamble for Big Data and Management Decision-Making slides.

\usetheme{Madrid}
\usecolortheme{default}
\usefonttheme{professionalfonts}

\usepackage{amsmath,amssymb,mathtools}
\usepackage{booktabs,array}
\usepackage{tikz}
\usepackage{graphicx}
\usepackage{listings}
\usepackage{xcolor}
\usetikzlibrary{arrows.meta,positioning,shapes.geometric}

\newcommand{\BDMFandolPath}{/usr/local/texlive/2025/texmf-dist/fonts/opentype/public/fandol/}
\setCJKmainfont[
  Path=\BDMFandolPath,
  UprightFont=FandolSong-Regular.otf,
  BoldFont=FandolSong-Regular.otf,
  ItalicFont=FandolKai-Regular.otf,
  BoldItalicFont=FandolKai-Regular.otf
]{FandolSong-Regular.otf}
\setCJKsansfont[
  Path=\BDMFandolPath,
  UprightFont=FandolSong-Regular.otf,
  BoldFont=FandolSong-Regular.otf
]{FandolSong-Regular.otf}
\setCJKmonofont[
  Path=\BDMFandolPath,
  UprightFont=FandolSong-Regular.otf,
  BoldFont=FandolSong-Regular.otf
]{FandolSong-Regular.otf}
\setmonofont{Menlo}

\definecolor{AcademicBlue}{RGB}{22,44,73}
\definecolor{AcademicTeal}{RGB}{22,119,119}
\definecolor{AcademicGray}{RGB}{76,84,95}
\definecolor{AcademicLight}{RGB}{245,247,250}
\definecolor{AcademicAmber}{RGB}{181,111,35}
\definecolor{CodeBg}{RGB}{248,248,248}

\setbeamercolor{structure}{fg=AcademicBlue}
\setbeamercolor{frametitle}{fg=white,bg=AcademicBlue}
\setbeamercolor{title}{fg=white,bg=AcademicBlue}
\setbeamercolor{block title}{fg=white,bg=AcademicBlue}
\setbeamercolor{block body}{bg=AcademicLight,fg=black}
\setbeamertemplate{navigation symbols}{}
\setbeamertemplate{footline}[frame number]
\setbeamertemplate{itemize item}{\raisebox{0.2ex}{\(\blacktriangleright\)}}
\setbeamerfont{title}{series=\mdseries}
\setbeamerfont{subtitle}{series=\mdseries}
\setbeamerfont{frametitle}{series=\mdseries}
\setbeamerfont{block title}{series=\mdseries}
\setbeamerfont{section in toc}{series=\mdseries}

\lstdefinestyle{pythonstyle}{
  basicstyle=\ttfamily\scriptsize,
  backgroundcolor=\color{CodeBg},
  keywordstyle=\color{AcademicBlue},
  commentstyle=\color{AcademicGray},
  stringstyle=\color{AcademicAmber},
  showstringspaces=false,
  breaklines=true,
  frame=single,
  rulecolor=\color{black!20},
  columns=fullflexible
}

\lstdefinestyle{sqlstyle}{
  language=SQL,
  basicstyle=\ttfamily\scriptsize,
  backgroundcolor=\color{CodeBg},
  keywordstyle=\color{AcademicBlue},
  commentstyle=\color{AcademicGray},
  stringstyle=\color{AcademicAmber},
  showstringspaces=false,
  breaklines=true,
  frame=single,
  rulecolor=\color{black!20},
  columns=fullflexible
}

\lstset{style=pythonstyle}

\hypersetup{
  colorlinks=true,
  linkcolor=AcademicBlue,
  urlcolor=AcademicTeal,
  pdfauthor={大数据与管理决策基础},
  pdfsubject={大数据与管理决策基础课程课件}
}


\title[大数据与管理决策基础]{第1讲：大数据与管理决策导论}
\subtitle{从经营症状到数据、模型、行动与反馈}
\author{《大数据与管理决策基础》}
\date{}

\begin{document}

\begin{frame}
  \titlepage
  \vspace{-0.4cm}
  \begin{center}
    \small
    核心命题：数据分析的终点不是模型准确率，而是可执行、可评估、可反馈的管理行动。
  \end{center}
\end{frame}

\begin{frame}{课程定位：方法论与技术桥梁}
\begin{center}
\resizebox{0.98\linewidth}{!}{%
\begin{tikzpicture}[
  node distance=0.85cm,
  box/.style={draw=AcademicBlue!45, fill=AcademicLight, rounded corners=1pt,
              minimum width=2.15cm, minimum height=0.9cm, align=center, font=\small},
  active/.style={draw=AcademicBlue, fill=AcademicBlue, text=white,
              minimum width=2.35cm, minimum height=1.0cm, align=center, font=\small},
  arrow/.style={-{Stealth[length=2mm]}, thick, AcademicTeal}
]
\node[box] (c1) {数智化\\运营概论};
\node[active, right=of c1] (c2) {大数据与\\管理决策基础};
\node[box, right=of c2] (c3) {企业运营数据\\挖掘与分析};
\node[box, right=of c3] (c4) {流程优化与\\数理决策};
\node[box, right=of c4] (c5) {综合实训};
\draw[arrow] (c1) -- (c2);
\draw[arrow] (c2) -- (c3);
\draw[arrow] (c3) -- (c4);
\draw[arrow] (c4) -- (c5);
\end{tikzpicture}
}
\end{center}

\vspace{0.4cm}
\begin{block}{本课程承担的角色}
将企业经营数据、统计建模、因果判断、优化决策与系统部署放入同一条教学主线。
\end{block}
\end{frame}

\begin{frame}{从经营症状开始，而不是从算法开始}
\begin{table}
\small
\centering
\begin{tabular}{p{0.18\linewidth}p{0.34\linewidth}p{0.34\linewidth}}
\toprule
经营症状 & 管理追问 & 典型指标\\
\midrule
库存积压 & 需求下降、补货过量还是渠道错配？ & 库存周转率、缺货率、滞销金额\\
交付延期 & 供应商、产能还是物流造成瓶颈？ & 准时交付率、平均延期天数\\
客户流失 & 哪些客户会流失，干预是否有效？ & 流失率、复购率、客户终身价值\\
产能瓶颈 & 哪个资源限制了系统吞吐量？ & 吞吐时间、资源利用率、等待时间\\
\bottomrule
\end{tabular}
\end{table}

\begin{alertblock}{教学要求}
任何分析任务都必须说明：谁会依据分析结果，对哪个业务对象，采取什么行动。
\end{alertblock}
\end{frame}

\begin{frame}{数据驱动管理决策的基本闭环}
\begin{center}
\resizebox{0.98\linewidth}{!}{%
\begin{tikzpicture}[
  stage/.style={draw=AcademicBlue!50, fill=AcademicLight, minimum width=1.65cm,
                minimum height=0.75cm, align=center, font=\small},
  action/.style={draw=AcademicAmber, fill=AcademicAmber!18, minimum width=1.65cm,
                minimum height=0.75cm, align=center, font=\small},
  arrow/.style={-{Stealth[length=2mm]}, thick, AcademicTeal}
]
\node[stage] (s1) {业务\\症状};
\node[stage, right=0.65cm of s1] (s2) {管理\\问题};
\node[stage, right=0.65cm of s2] (s3) {数据\\表达};
\node[stage, right=0.65cm of s3] (s4) {指标\\体系};
\node[stage, right=0.65cm of s4] (s5) {分析\\模型};
\node[action, right=0.65cm of s5] (s6) {管理\\行动};
\node[stage, right=0.65cm of s6] (s7) {反馈\\评价};
\foreach \i/\j in {s1/s2,s2/s3,s3/s4,s4/s5,s5/s6,s6/s7}
  \draw[arrow] (\i) -- (\j);
\end{tikzpicture}
}
\end{center}

\vspace{0.4cm}
\[
\text{业务症状}\rightarrow \text{数据表达}\rightarrow \text{模型判断}\rightarrow
\text{行动选择}\rightarrow \text{反馈治理}.
\]
\end{frame}

\begin{frame}{五类分析回答不同问题}
\begin{table}
\small
\centering
\begin{tabular}{p{0.15\linewidth}p{0.25\linewidth}p{0.46\linewidth}}
\toprule
类型 & 核心问题 & 示例\\
\midrule
描述分析 & 发生了什么？ & 本月销售额、订单量、库存水平。\\
诊断分析 & 为什么发生？ & 哪些区域、产品、供应商造成延期。\\
预测分析 & 将会发生什么？ & 哪些订单可能延期，哪些客户可能流失。\\
因果分析 & 行动是否有效？ & 优惠券是否真正降低流失率。\\
处方分析 & 应该怎么做？ & 在预算约束下给哪些客户发券。\\
\bottomrule
\end{tabular}
\end{table}

\begin{block}{边界意识}
预测分数不是行动建议；相关关系不是因果效应；仪表盘不是完整决策系统。
\end{block}
\end{frame}

\begin{frame}{预测、因果与行动的分工}
\begin{table}
\small
\centering
\begin{tabular}{p{0.24\linewidth}p{0.31\linewidth}p{0.31\linewidth}}
\toprule
问题 & 数学对象 & 管理含义\\
\midrule
预测 & \(P(Y=1\mid X=x)\) & 哪些对象风险更高。\\
因果 & \(\mathbb E[Y(a)-Y(a')\mid X=x]\) & 行动是否真的改变结果。\\
处方 & \(a^*(x)\) & 在成本和约束下选择行动。\\
\bottomrule
\end{tabular}
\end{table}

\begin{alertblock}{课堂提醒}
高风险对象不一定最值得行动；行动优先级还取决于行动效果、对象价值、资源约束和可审计责任。
\end{alertblock}
\end{frame}

\begin{frame}{管理决策的统一数学表达}
给定企业状态 \(X=x\)，在可行行动集合 \(\mathcal A(x)\) 中选择行动：
\[
a^*(x)
=
\arg\min_{a\in\mathcal A(x)}
\mathbb E\!\left[L(Y,a)\mid X=x\right]
+
\lambda C(a).
\]

\begin{table}
\small
\centering
\begin{tabular}{cl}
\toprule
符号 & 含义\\
\midrule
\(X\) & 当前企业状态数据：订单、库存、客户、产能、流程日志\\
\(Y\) & 未来经营结果：延期、流失、需求、利润\\
\(a\) & 管理行动：加急、补货、发券、重排产能\\
\(L(Y,a)\) & 行动后产生或避免的业务损失\\
\(C(a)\) & 行动成本、资源占用或约束惩罚\\
\bottomrule
\end{tabular}
\end{table}
\end{frame}

\begin{frame}{CRISP-DM 是起点，现代企业还需要治理}
\begin{columns}[T]
\begin{column}{0.48\linewidth}
\begin{block}{CRISP-DM}
\begin{enumerate}
  \item Business Understanding
  \item Data Understanding
  \item Data Preparation
  \item Modeling
  \item Evaluation
  \item Deployment
\end{enumerate}
\end{block}
\end{column}
\begin{column}{0.48\linewidth}
\begin{block}{现代闭环补充}
\begin{itemize}
  \item 数据质量与指标治理；
  \item 元数据、血缘与可复现环境；
  \item 模型监控、漂移检测与审计；
  \item 权限控制与 human-in-the-loop；
  \item 行动日志与反馈回流。
\end{itemize}
\end{block}
\end{column}
\end{columns}
\end{frame}

\begin{frame}{案例：制造企业交付延期风险预警}
\begin{table}
\small
\centering
\begin{tabular}{p{0.18\linewidth}p{0.68\linewidth}}
\toprule
层次 & 内容\\
\midrule
业务症状 & 准时交付率下降，关键客户投诉增加。\\
管理问题 & 哪些订单需要提前加急或重排产能？\\
候选数据 & 订单、供应商、库存、生产排程、物流轨迹、流程日志。\\
分析模型 & 延期风险评分、瓶颈诊断、产能分配优化。\\
管理行动 & 加急采购、产能重排、替代物流或客户沟通。\\
反馈指标 & 准时交付率、加急成本、客户满意度。\\
\bottomrule
\end{tabular}
\end{table}

\[
\text{风险预测}\rightarrow\text{行动推荐}\rightarrow\text{情景比较}\rightarrow\text{行动日志}\rightarrow\text{反馈评价}.
\]
\end{frame}

\begin{frame}[fragile]{配套代码：问题框架的结构化表达}
\begin{lstlisting}[language=Python]
from dataclasses import dataclass

@dataclass(frozen=True)
class OperatingSymptom:
    name: str
    business_question: str
    candidate_data: list[str]
    kpis: list[str]
    analysis_type: str
    decision_action: str
\end{lstlisting}

\begin{block}{运行方式}
\begin{lstlisting}[language=bash]
PYTHONPATH=src python3 src/bdm_decision/cases/week01_problem_framing.py
\end{lstlisting}
\end{block}
\end{frame}

\begin{frame}{课堂练习：15 分钟完成问题拆解}
\begin{enumerate}
  \item 选择一个业务症状：库存、交付、客户、产能或自选案例；
  \item 写出一个包含行动对象的管理问题；
  \item 列出至少四类数据表或事件日志；
  \item 定义至少五个 KPI，并说明管理含义；
  \item 判断涉及描述、诊断、预测、因果或处方分析中的哪些类型；
  \item 提出至少两个可执行管理行动；
  \item 设计行动后的反馈指标。
\end{enumerate}
\end{frame}

\begin{frame}{本周交付与下周衔接}
\begin{block}{本周交付}
\begin{itemize}
  \item 讲义：\texttt{docs/lecture\_notes/01\_intro.tex}
  \item 课件：\texttt{slides/01\_intro/week01\_data\_driven\_decision\_intro\_beamer.tex}
  \item 代码：\texttt{src/bdm\_decision/cases/week01\_problem\_framing.py}
  \item 作业：\texttt{assignments/assignment\_01\_data\_modeling.md}
\end{itemize}
\end{block}

\begin{alertblock}{下周主题}
企业数据与业务对象：从关系表、实体、事件与状态出发，建立订单--客户--产品--库存数据模型。
\end{alertblock}
\end{frame}

\end{document}
